% !PRINTED 256
\weberpage{256}
Sondert man in (15) die erste Potenz von $x$ ab und erhebt ins
Quadrat, so erhält man durch Einsetzen von (21) die gesuchte
Multiplikatorgleichung für den 9ten Transformationsgrad. Sie erhält
eine einfachere Gestalt, wenn man
\weq{22}{(y+3)^3 = t}
setzt:
\weq{23}{[j(\omega) - 27\cdot64](t-27) = (t^2 - 36t + 27\cdot8)^2}
oder
\weq{24}{j(\omega)(t-27) = t(t-24)^3.}

Ganz ähnlich kann man beim 25sten Transformationsgrad verfahren.

Wenn wir
\weq{25}{x = \left(\frac{\eta\left(\frac{\omega}{5}\right)}{\eta(\omega)}\right)^2, \quad y = \frac{\eta\left(\frac{\omega}{25}\right)}{\eta(\omega)}}
setzen, so haben wir zunächst nach 4. (8):
\weq{26}{x^6 + 10x^3 - \gamma_2(\omega)x + 5 = 0,}
woraus, wie oben, die beiden Gleichungen
\[
5^5 + 10\cdot5^2x^3 - \gamma_2\left(\frac{\omega}{5}\right)x^5 + x^6 = 0,
\]
\[
y^{12} + 10y^6x^3 - \gamma_2\left(\frac{\omega}{5}\right)y^2x^5 + 5x^6 = 0,
\]
und durch Elimination von $\gamma_2\left(\frac{\omega}{5}\right)$
\[
x^6 - 10x^3y^2(5+y^2) = y^2(y^8 + 5y^6 + 5^2y^4 + 5^3y^2 + 5^4).
\]
Diese Gleichung nach $x^3$ aufgelöst, ergibt:
\weq{27}{x^3 = y^5 + 5y^4 + 15y^3 + 25y^2 + 25y,}
und wenn wir zur Abkürzung die rechte Seite dieser Gleichung mit
$\chi(y)$ bezeichnen, nach (26)
\weq{28}{j(\omega)\chi(y) = [\chi(y)^2 + 10\chi(y) + 5]^3.}
Dies ist die gesuchte Gleichung 30sten Grades für $y$.

\begin{center}
§ 73.\quad \gesperrt{Die Schlaeflischen Modulargleichungen.}
\end{center}

Zu einfacheren Transformationsgleichungen gelangt man, wenn man dem
Rationalitätsbereich, der bis jetzt aus den rationalen Funktionen von
$j(\omega)$ bestand, die Größe $f(\omega)^{24}$ adjungiert. Diesem
%CONT>>
